\SourceLeaf{079}
\[
\xi=a,\quad\eta=b,\quad\zeta=c\quad\text{etc.}
\]
valores ipsarum \(x,y,z\) etc. hi
\[
x=\alpha,\quad y=\beta,\quad z=\gamma\quad\text{etc.}
\]
ita ut habeatur
\[
\begin{aligned}
\alpha&=a[\alpha\alpha]+b[\alpha\beta]+c[\alpha\gamma]+\text{etc.}\\
\beta&=a[\alpha\beta]+b[\beta\beta]+c[\beta\gamma]+\text{etc.}\\
\gamma&=a[\alpha\gamma]+b[\beta\gamma]+c[\gamma\gamma]+\text{etc.}
\end{aligned}
\]
etc. Perinde valoribus
\[
\xi=a',\quad\eta=b',\quad\zeta=c'\quad\text{etc.}
\]
respondere supponemus hos
\[
x=\alpha',\quad y=\beta',\quad z=\gamma'\quad\text{etc.}
\]
nec non his
\[
\xi=a'',\quad\eta=b'',\quad\zeta=c''\quad\text{etc.}
\]
sequentes
\[
x=\alpha'',\quad y=\beta'',\quad z=\gamma''\quad\text{etc.}
\]
et sic porro.

His positis combinatio aequationum (4), (9) suppeditat
\[
\begin{aligned}
A&=\alpha e+\alpha'e'+\alpha''e''+\text{etc.}\\
B&=\beta e+\beta'e'+\beta''e''+\text{etc.}\\
C&=\gamma e+\gamma'e'+\gamma''e''+\text{etc.}
\end{aligned}
\]
etc. Quare quum habeatur \(S=\mathfrak A A+\mathfrak B B+\mathfrak C C+\text{etc.}\), patet fieri
\[
\begin{aligned}
S={}&(ae+a'e'+a''e''+\text{etc.})(\alpha e+\alpha'e'+\alpha''e''+\text{etc.})\\
&+(be+b'e'+b''e''+\text{etc.})(\beta e+\beta'e'+\beta''e''+\text{etc.})\\
&+(ce+c'e'+c''e''+\text{etc.})(\gamma e+\gamma'e'+\gamma''e''+\text{etc.})+\text{etc.}
\end{aligned}
\]
\Article{16}
Institutionem observationum, per quas valores quantitatum \(v,v',v''\) etc. erroribus fortuitis \(e,e',e''\) etc. affectos obtinemus, considerare possumus tamquam
\SourceLeaf{080}
experimentum, quod quidem singulorum errorum commissorum magnitudinem docere non valet, attamen, praeceptis quae supra explicavimus adhibitis, valorem quantitatis \(S\) subministrat, qui per formulam modo inventam est functio data illorum errorum. In tali experimento errores fortuiti utique alii maiores alii minores prodire possunt; sed quo plures errores concurrunt, eo maior spes aderit, valorem quantitatis \(S\) in experimento singulari a valore suo medio parum deviaturum esse. Rei cardo itaque in eo vertitur, ut valorem medium quantitatis \(S\) stabiliamus. Per principia in commentatione priore exposita, quae hic repetere superfluum esset, invenimus hunc valorem medium
\[
\begin{aligned}
&=(a\alpha+b\beta+c\gamma+\text{etc.})mm\\&\qquad+(a'\alpha'+b'\beta'+c'\gamma'+\text{etc.})m'm'\\
&\qquad+(a''\alpha''+b''\beta''+c''\gamma''+\text{etc.})m''m''+\text{etc.}
\end{aligned}
\]
Denotando errorem medium observationum talium, quarum pondus \(=1\), per \(\mu\), ita ut sit \(\mu\mu=pmm=p'm'm'=p''m''m''\) etc., expressio modo inventa ita exhiberi potest:
\[
\begin{aligned}
&\left(\frac{a\alpha}{p}+\frac{a'\alpha'}{p'}+\frac{a''\alpha''}{p''}\ \text{etc.}\right)\mu\mu
\\&\qquad+\left(\frac{b\beta}{p}+\frac{b'\beta'}{p'}+\frac{b''\beta''}{p''}+\text{etc.}\right)\mu\mu\\
&\qquad+\left(\frac{c\gamma}{p}+\frac{c'\gamma'}{p'}+\frac{c''\gamma''}{p''}+\text{etc.}\right)\mu\mu+\text{etc.}
\end{aligned}
\]
Sed aggregatum \(\frac{a\alpha}{p}+\frac{a'\alpha'}{p'}+\frac{a''\alpha''}{p''}+\text{etc.}\) invenitur
\[
=[aa].[\alpha\alpha]+[ab].[\alpha\beta]+[ac].[\alpha\gamma]+\text{etc.}
\]
adeoque \(=1\), uti e nexu aequationum (6), (7) facile intelligitur. Perinde fit
\[
\begin{aligned}
\frac{b\beta}{p}+\frac{b'\beta'}{p'}+\frac{b''\beta''}{p''}+\text{etc.}&=1\\[5pt]
\frac{c\gamma}{p}+\frac{c'\gamma'}{p'}+\frac{c''\gamma''}{p''}+\text{etc.}&=1
\end{aligned}
\]
et sic porro.

Hinc tandem valor medius ipsius \(S\) fit \(=\sigma\mu\mu\), quatenusque igitur valorem fortuitum ipsius \(S\) pro medio adoptare licet, erit \(\mu=\BareRoot{\frac S\sigma}\).

\Article{17}
Quanta fides huic determinationi habenda sit, diiudicare oportet per errorem medium vel in ipsa vel in ipsius quadrato metuendum: posterior erit radix
\SourceLeaf{081}
quadrata valoris medii expressionis
\[
\left(\frac S\sigma-\mu\mu\right)^2
\]
cuius evolutio absolvetur per ratiocinia similia iis, quae in commentatione priore artt. 39 sqq. exposita sunt. Quibus brevitatis caussa hic suppressis, formulam ipsam tantum hic apponimus. Scilicet error medius in determinatione quadrati \(\mu\mu\) metuendus exprimitur per
\[
\BareRoot{\left(\frac{2\mu^4}{\sigma}+\frac{\nu^4-3\mu^4}{\sigma\sigma}\cdot N\right)}
\]
denotante \(\nu^4\) valorem medium biquadratorum errorum. quorum pondus \(=1.\) atque \(N\) aggregatum
\[
\begin{aligned}
&(a\alpha+b\beta+c\gamma+\text{etc.})^2+(a'\alpha'+b'\beta'+c'\gamma'+\text{etc.})^2\\
&\qquad+(a''\alpha''+b''\beta''+c''\gamma''+\text{etc.})^2\text{etc.}
\end{aligned}
\]
Hoc aggregatum in genere ad formam simpliciorem reduci nequit, sed simili modo ut in art. 40 prioris commentationis ostendi potest, eius valorem semper contineri intra limites \(\pi\) et \(\frac{\sigma\sigma}{\pi}\). In hypothesi ea, cui theoria quadratorum minimorum ab initio superstructa erat, terminus hoc aggregatum continens, propter \(\nu^4=3\mu^4\), omnino excidit, praecisioque, quae errori medio, per formulam \(\BareRoot{\frac S\sigma}\) determinato, tribuenda est, eadem erit, ac si ex \(\sigma\) erroribus exacte cognitis secundum artt. 15,16 prioris commentationis erutus fuisset.

\Article{18}
Ad compensationem observationum duo, ut supra vidimus, requiruntur: primum, ut aequationum conditionalium correlata, i. e. numeri \(A,B,C\) etc. aequationibus (12) satisfacientes eruantur, secundum, ut hi numeri in aequationibus (10) substituantur. Compensatio hoc modo prodiens dici poterit \textit{perfecta} seu \textit{completa}, ut distinguatur a compensatione \textit{imperfecta} seu \textit{manca}: hac scilicet denominatione designabimus, quae resultant ex iisdem quidem aequationibus (10), sed substratis valoribus quantitatum \(A,B,C\) etc., qui non satisfaciunt aequationibus (12), i. e. qui vel parti tantum satisfaciunt vel nullis. Quod vero attinet ad tales observationum mutationes, quae sub formulis (10) comprehendi nequeunt, a disquisitione praesenti, nec non a denominatione compensationum exclusae sunto. Quum, quatenus aequationes (10) locum habent, aequationes (13) ipsis (12) omnino sint aequivalentes, illud discrimen ita quoque enunciari potest: Ob%
